\section{Conclusions}
\label{sec:conclusions}
\begin{enumerate}
  \item The split-spectrum screens of NISAR's PR GUNWs copy the coseismic deformation near the 2026 Panama and
        Venezuela ruptures with the opposite sign. Within 10\,km of the trace, the screen's high-pass regresses on
        the data's at $-0.02$ to $-0.12$ ($r$ down to $-0.82$). On Panama's D141 pair the screen also steps 0.43\,m
        across the trace.
  \item On identical cells, with an identical filter and main-band unwrapping, two mechanisms account for most of it.
        \begin{itemize}
          \item The whole-cycle corrections read the deformation's curvature as unwrapping errors, and copy
                4.5--4.8\% of the high-passed deformation at any slip.
          \item Unwrapping a field that carries deformation (the side band, or the double difference) fails across a
                rupture that splits the scene: by one cycle of the double difference (1.4\,m of LOS) on D141.
        \end{itemize}
        On D141 the product's own main band, 25--32 cycles off across the rupture, carries most of its copy. On the
        Venezuela tracks, the corrections on matched cells reproduce 25--52\% of the products' leak beyond ours, and
        neither the main band nor the bands' footprint mismatch explains the rest.
  \item Removing the scaled, unwrapped main band from the wrapped side band before anything is unwrapped gives the
        same estimate as the reformulated split spectrum and the GUNW's, with nothing in the unwrapped field for an
        unwrapper or a correction to mistake. On every synthetic its error near the rupture stays at its zero-slip
        level. It needs the main band unwrapped correctly, as the deformation does, and a cut-off block's tie carries
        into the screen at 3.21\,rad per cycle.
  \item The same 3.21\,rad per cycle ties a block that a cut along the rupture isolates. The raw estimate's step across
        the gap reads the block's integer wherever the ionosphere is smooth over 10--20\,km: the median over 20\,km
        windows read no wrong integer on 565 fake boundaries in eight scenes, and 0.02 across a Venezuela rupture whose
        integer is known. On three Panama pairs it finds the pipelines' range-offset ties a cycle short. Read with the same
        estimator, the offsets agree with it: the pipelines' far-field planes were at fault, not the offsets. Ionospheric waves, and
        local cycle errors in a third of the cells beside the gap, bias it by a few tenths of a cycle; two estimators
        that fail differently, a polynomial and kriging, check each other.
  \item What remains near the fault on the synthetics, the same with and without the rupture, is the cells' handling of
        steep fringes, the fill's width and the lag of ISCE3's weighted mean across the damage zone. Taking the main
        band out of every pixel, matching the fill to the gap and making the filter's final step a weighted local
        plane bring it from 16 to 2\,mm within 2\,km of the trace. On real data these are millimetre corrections.
        Next to the D141 rupture the screen is limited by the main band's local integers, which the narrower recipes
        follow and the production screen's wide fill averages away.
  \item On two Antarctic pairs the operational screen is ours plus whole cycles:
        \begin{itemize}
          \item the main band's split between its components, 9 and 8 cycles;
          \item one cycle of the double difference over 12\,000\,km$^2$ beside a glacier's shear margins;
          \item about two main-band cycles from the correction passes over a fast trunk.
        \end{itemize}
        Where the screens differ beyond those integers, the range offsets side with ours. They also tie the components as
        the side band does across a 160\,rad step, where velocity maps read one to two cycles low.
  \item For the operational processor the change is one line: multiply the double difference by
        $\exp(i\,\Delta f\,\phiA/\fA)$ before unwrapping, and divide it out after. The rupture still needs a cut, and
        the side band then ties the blocks.
\end{enumerate}
